\documentclass[11pt, a4paper]{article}
\usepackage[top=1in, bottom=1in, left=1.1in, right=1.1in]{geometry}
\usepackage{lmodern}
\usepackage{amsmath, amssymb}
\usepackage{booktabs}
\usepackage{enumitem}
\usepackage{xcolor}
\usepackage[colorlinks=true,linkcolor=black,citecolor=black,urlcolor=blue]{hyperref}

\makeatletter
\renewcommand{\section}{\@startsection{section}{1}{\z@}%
  {10pt plus2pt minus2pt}{4pt}{\large\bfseries}}
\renewcommand{\subsection}{\@startsection{subsection}{2}{\z@}%
  {7pt plus1pt minus1pt}{2pt}{\normalsize\bfseries}}
\makeatother
\setlength{\parskip}{4pt}
\setlength{\parindent}{0pt}
\setlist{nosep, leftmargin=1.4em}

\begin{document}

\begin{center}
{\Large\bfseries Alpha Fit --- Purpose \& Transition Note}\\[4pt]
{\small VSEPR-Sim Atomistic Engine \quad|\quad Polarizability Module}
\end{center}
\vspace{-0.4em}\hrule\vspace{0.6em}

% ─────────────────────────────────────────────────────────────────────────────
\section{What Alpha Fitting Was}

The \textit{alpha fitting} subsystem
(\texttt{tools/fit\_alpha\_model.cpp}, \texttt{tools/auto\_fit\_alpha.cpp})
was an offline coordinate-descent trainer for the scalar isotropic
polarizability model, referred to internally as \textbf{Alpha Method D}.

The model estimated atomic polarizability $\alpha(Z)$ in \AA$^3$ for
elements $Z > 118$ (i.e.\ beyond the empirical table) using the formula:

\[
\alpha(Z) \;=\; k_p \cdot c_b \cdot r_\text{eff}^3 \cdot g_\text{bind}
             \cdot g_f \;+\; \text{blob}
\]

\noindent where:

\begin{itemize}
  \item $k_p$ --- per-period scaling ($k_\text{period}[1\ldots7]$)
  \item $c_b$ --- per-block correction ($c_\text{block}[s,p,d,f]$)
  \item $r_\text{eff} = r_\text{cov}\,(1 + c_\text{rel}\,Z^2)\,
        (1 - \beta_f\,n_f/14)$ --- effective covalent radius with
        relativistic and lanthanide-contraction corrections
  \item $g_\text{bind} = 1/(1 + b_\text{bind}\,I_1)$ --- ionization
        stiffness gate ($I_1$ = first ionisation energy, eV)
  \item $g_f$ --- f-block multiplicative correction using
        $n_f$, half-shell bump ($a_{f2}$), and full-shell correction ($a_{f3}$)
  \item $\text{blob}$ --- additive f-block anomaly:
        $\text{blob} = b_\text{lin}\,n_f + b_\text{half}\,\delta_{n_f=7}
         + b_\text{full}\,\delta_{n_f=14}$
\end{itemize}

The trainer minimised a weighted Huber loss in log-space over a reference CSV
(\texttt{data/polarizability\_ref.csv}) covering Z\,=\,1--118, with
chemistry-aware weights (noble gases $\times 4$, halogens $\times 2.5$,
early actinides $\times 2$, relativistic Hg $\times 2.5$) and a
group-wise monotonicity smoothness penalty.

% ─────────────────────────────────────────────────────────────────────────────
\section{Why It Is a Dead End}

\begin{enumerate}
  \item \textbf{Zero impact on Z\,=\,1--118.}\;
        \texttt{alpha\_predict()} returns the empirical table directly for
        all known elements; the fitted model is only invoked for $Z > 118$.
        No physical system simulated by the engine today uses super-heavy elements.

  \item \textbf{Circular validation.}\;
        The reference data and the fitter are both internal; there is no
        experimental observable in the production simulation loop that
        depends on $\alpha(Z > 118)$.

  \item \textbf{Parameter drift.}\;
        Successive alpha model versions changed the field layout
        (\texttt{c\_shell}, \texttt{a\_rel}, \texttt{c\_orb} $\to$
        \texttt{beta\_f}, \texttt{c\_rel}, \texttt{b\_bind}, etc.),
        causing repeated stale-reference build failures in consumer tools
        without any improvement to production accuracy.

  \item \textbf{Scope mismatch.}\;
        The engine's validated domain is $Z = 1$--102 (periodic table
        through Nobelium).  Fitting a model for $Z > 118$ is pre-emptive
        infrastructure for physics that does not yet exist in the engine.
\end{enumerate}

% ─────────────────────────────────────────────────────────────────────────────
\section{Transition Decision}

The alpha fitting tools are \textbf{retained in the repository as
reference} but are no longer on the critical build path.  The transition is:

\begin{center}
\renewcommand{\arraystretch}{1.3}
\begin{tabular}{ll}
\toprule
\textbf{Old approach} & \textbf{New approach} \\
\midrule
Fitted generative model for $Z > 118$ & Return default value (1.76~\AA$^3$) \\
Offline trainer binaries in build & Tools excluded from default targets \\
Parameter struct fields as build API & Model struct treated as opaque \\
Validation via training loss & Validation via production test suite \\
\bottomrule
\end{tabular}
\end{center}

The freed build budget and code-review attention are redirected to the
formation pipeline (SS6), the heat-gated reaction control (SS8b), and
the 35-test validation campaign (SS12).

% ─────────────────────────────────────────────────────────────────────────────
\section{What Remains Valid}

\begin{itemize}
  \item \texttt{atomistic/models/alpha\_model.hpp} ---
        the \texttt{AlphaModelParams} struct and \texttt{alpha\_predict()}
        function remain unchanged.  Callers using Z\,=\,1--118 are
        unaffected.
  \item \texttt{data/polarizability\_ref.csv} ---
        the reference dataset is kept as documentation of the empirical values.
  \item The Thole damping model, SCF polarization loop, and QEq charge
        equilibration (SS8) are independent of the fitter and continue
        to be developed.
\end{itemize}

\vspace{1em}
\hrule
\vspace{0.4em}
{\small\textit{Status: Transitioned. Tools retained for reference.
Formation engine version 2.5.2.X and later.}}

\end{document}
